\scanpage{57}
\mspage{\handwritten{- 9 -}}
% archive sheet number: 56
 des séries entières $g_k$, et en substituant à $z_k$ dans $\handwritten{f}$ chacune des séries entières en $v_1,\ldotp\ldotp,v_r$ obtenues : $\handwritten{2^{\mathrm{o}}}$ soit en substituant $g_k$ à $z_k$ dans $\handwritten{f}$, et en substituant $\handwritten{\unsure{\gamma|y}}_j$ à $u_j$ dans la série entière en $u_1,\ldots,u_q$ obtenue.

\uline{\handwritten{4}}. \uline{La notion de fonction holomorphe}.

Le th.2 entraîne en particulier la proposition suivante :

PROPOSITION 4.- \uline{Soit} $\handwritten{f}(z_1,\ldots,z_p)=\handwritten{\sum}_{\handwritten{(n_k)}}\handwritten{\unsure{c|C}}_{n_1n_2\ldotp\ldotp n_p}z_1^{n_1}z_2^{n_2}\ldots z_p^{n_p}$ \uline{la somme d'une série entière sommable dans un polycylindre ouvert} $P$ \uline{de centre} $0$ \uline{défini par} $\handwritten{|}z_k\handwritten{|}\handwritten{<}r_k$ $(1\handwritten{\leq}k\handwritten{\leq}p)$. \uline{Soit} $(a_k)$ \uline{un point quelconque de} $P$ ; \uline{pour tout point} $(u_k)_{1\handwritten{\leq}k\handwritten{\leq}p}$ \uline{du polycylindre défini par} $\handwritten{|}u_k\handwritten{|}\handwritten{<}r_k-\handwritten{|}a_k\handwritten{|}$, \uline{les} $p$ \uline{polynomes} $a_k+u_k$ \uline{sont substituables aux} $z_k$ \uline{dans la série entière} $\handwritten{c}_{n_1\ldotp\ldotp n_p}z_1^{n_1}\ldotp\ldotp z_p^{n_p}$, \uline{et la somme de la série entière} (absolument sommable) \uline{par rapport aux} $u_k$ \uline{ainsi obtenue est égale à} $\handwritten{f}(a_1+u_1,\ldots,a_p+u_p)$.

En effet, avec les notations du th.2, on a ici $h_k(\handwritten{|}u_1\handwritten{|},\ldots,\handwritten{|}u_p\handwritten{|})==\handwritten{|}a_k\handwritten{|}+\handwritten{|}u_k\handwritten{|}\handwritten{<}r_k$ par hypothèse.

En d'autres termes, il existe un polycylindre ouvert de centre $(a_1,\ldots,a_p)$ dans lequel $\handwritten{f}(z_1,\ldots,z_p)$ est égale à la somme d'une série entière absolument sommable, par rapport à $z_k-a_k$ $(1\handwritten{\leq}k\handwritten{\leq}p)$.

DÉFINITION 2.- \uline{Soit} $A$ \uline{un ensemble ouvert dans} ${K}^{p}$. \uline{On dit qu'une application} $\handwritten{f}$ \uline{de} $A$ \uline{dans} $E$ \uline{est une fonction holomorphe dans} $A$ \uline{si}, \uline{pour tout point} $\handwritten{a}=(a_k)_{1\handwritten{\leq}k\handwritten{\leq}p}$ \uline{de} $A$, \uline{il existe un polycylindre ouvert de centre} $\handwritten{a}$, \uline{contenu dans} $A$, \uline{et dans lequel} $\handwritten{f}(z_1,z_2,\ldots,z_p)$ \uline{soit égale à la somme d'une série entière absolument sommable}, \uline{par rapport aux variables} $z_k-a_k$ $(1\handwritten{\leq}k\handwritten{\leq}p)$.

La prop.4 montre donc que :

PROPOSITION 5.- \uline{Si une série entière par rapport aux variables} $z_k$ $(1\handwritten{\leq}k\handwritten{\leq}p)$ \uline{est sommable dans un polycylindre ouvert de centre} $0$,
